% Dynamic Pareto Filtering Section
% This section describes the context-dependent Pareto frontier pruning mechanism

\subsection{Dynamic Pareto Filtering}
\label{sec:dynamic_pareto_filtering}

\subsubsection{Motivation}
In multi-model routing with heterogeneous cost-quality profiles, the naive approach of evaluating all $K$ models at each time step incurs unnecessary computational overhead and exploration risk. Many models are \textit{strictly dominated} for a given context: a model $m$ is dominated if there exists another model $m'$ that is both cheaper ($c_{m'} < c_m$) and predicted to perform better ($\hat{r}_{m'}(x) > \hat{r}_m(x)$) on the current prompt $x$.

We introduce \textbf{Dynamic Pareto Filtering}, a context-aware pruning mechanism that restricts the LinUCB selection policy to the local Pareto frontier, eliminating dominated models before utility maximization.

\subsubsection{Methodology}

For each incoming prompt $x$ with feature representation $\mathbf{x} \in \mathbb{R}^d$, we compute contextual quality predictions for all models:
\begin{equation}
    \hat{r}_m(x) = \mathbf{x}^\top \hat{\boldsymbol{\theta}}_m = \mathbf{x}^\top \mathbf{A}_m^{-1} \mathbf{b}_m
\end{equation}

where $\mathbf{A}_m$ and $\mathbf{b}_m$ are the sufficient statistics maintained by the LinUCB algorithm.

\paragraph{Pareto Dominance Relation.}
Model $m'$ \textit{Pareto-dominates} model $m$ for context $x$ if:
\begin{equation}
    \hat{r}_{m'}(x) \geq \hat{r}_m(x) \quad \text{and} \quad c_{m'} \leq c_m \quad \text{and} \quad (\hat{r}_{m'}(x) > \hat{r}_m(x) \text{ or } c_{m'} < c_m)
\end{equation}

The context-dependent Pareto frontier is defined as:
\begin{equation}
    \mathcal{P}(x) = \left\{ m \in \mathcal{M} : \nexists m' \in \mathcal{M} \text{ such that } m' \text{ dominates } m \text{ for } x \right\}
\end{equation}

\paragraph{Filtered Selection Policy.}
The router's selection policy becomes:
\begin{equation}
    m^* = \argmax_{m \in \mathcal{P}(x)} \left[ w_q \cdot \text{UCB}_m(x) - w_c \cdot c_m \right]
\end{equation}

where:
\begin{itemize}
    \item $\text{UCB}_m(x) = \hat{r}_m(x) + \alpha \sqrt{\mathbf{x}^\top \mathbf{A}_m^{-1} \mathbf{x}}$ is the Upper Confidence Bound
    \item $w_q, w_c \in [0,1]$ are quality and cost preference weights ($w_q + w_c = 1$)
    \item $c_m$ is the normalized cost of model $m$
\end{itemize}

\subsubsection{Computational Efficiency}

\paragraph{Complexity Analysis.}
\begin{itemize}
    \item \textbf{Naive Approach}: $O(K \cdot d^2)$ to compute UCB for all $K$ models (matrix-vector products $\mathbf{A}_m^{-1} \mathbf{x}$)
    \item \textbf{With Pareto Filtering}: $O(K \log K + |\mathcal{P}(x)| \cdot d^2)$ where $|\mathcal{P}(x)| \ll K$ typically
    \item \textbf{Speedup}: For $K=9$ models, empirical measurements show $|\mathcal{P}(x)| \approx 3$, yielding a $3\times$ reduction in UCB computations
\end{itemize}

The $O(K \log K)$ term comes from sorting models by cost to construct the Pareto frontier, which is negligible compared to the $O(d^2)$ matrix operations.

\subsubsection{Theoretical Properties}

\paragraph{Optimality Preservation.}
Dynamic Pareto Filtering does \textit{not} sacrifice optimality. If model $m^*$ is the true optimal choice under the full utility function, then $m^* \in \mathcal{P}(x)$ by definition (an optimal model cannot be dominated). Therefore:
\begin{equation}
    \argmax_{m \in \mathcal{M}} U_m(x) = \argmax_{m \in \mathcal{P}(x)} U_m(x)
\end{equation}

\paragraph{Regret Bound.}
Since Pareto filtering only removes strictly suboptimal actions, the LinUCB regret bound remains unchanged:
\begin{equation}
    \text{Regret}_T = O\left( d \sqrt{T \log T} \right)
\end{equation}

However, the \textit{constant factors} improve due to reduced exploration variance (fewer arms to explore).

\subsubsection{Practical Benefits}

\begin{enumerate}
    \item \textbf{Reduced Exploration Risk}: By excluding dominated models, we avoid wasting exploration budget on arms that are provably suboptimal for the current context.
    
    \item \textbf{Faster Convergence}: Concentrating updates on Pareto-optimal models accelerates learning of their true quality parameters.
    
    \item \textbf{Interpretability}: The Pareto frontier provides a natural explanation for why certain models were excluded (``Model X was both slower and predicted to be worse than Model Y'').
    
    \item \textbf{Dynamic Adaptation}: The frontier adapts to each context. For example:
    \begin{itemize}
        \item \textit{Easy prompts}: Cheap models dominate, $|\mathcal{P}(x)| \approx 2$ (one cheap model + one mid-tier fallback)
        \item \textit{Hard prompts}: Expensive models dominate, $|\mathcal{P}(x)| \approx 3$ (multiple frontier models compete)
    \end{itemize}
\end{enumerate}

\subsubsection{Implementation Details}

The Pareto frontier construction algorithm:

\begin{algorithm}[t]
\caption{Dynamic Pareto Filtering}
\label{alg:pareto_filtering}
\begin{algorithmic}[1]
\Require Context $\mathbf{x}$, Models $\mathcal{M} = \{m_1, \ldots, m_K\}$, Costs $\{c_1, \ldots, c_K\}$, LinUCB parameters $\{\mathbf{A}_m, \mathbf{b}_m\}$
\Ensure Pareto frontier $\mathcal{P}(x)$

\State Compute quality predictions: $\hat{r}_m \gets \mathbf{x}^\top \mathbf{A}_m^{-1} \mathbf{b}_m$ for all $m \in \mathcal{M}$
\State Sort models by cost: $\mathcal{M}_{sorted} \gets \text{sort}(\mathcal{M}, \text{key}=c_m)$
\State Initialize frontier: $\mathcal{P}(x) \gets \emptyset$
\State Initialize max quality: $r_{max} \gets -\infty$

\For{$m \in \mathcal{M}_{sorted}$} \Comment{Iterate from cheapest to most expensive}
    \If{$\hat{r}_m > r_{max}$} \Comment{Non-dominated: better quality than all cheaper models}
        \State $\mathcal{P}(x) \gets \mathcal{P}(x) \cup \{m\}$
        \State $r_{max} \gets \hat{r}_m$
    \EndIf
\EndFor

\State \Return $\mathcal{P}(x)$
\end{algorithmic}
\end{algorithm}

\paragraph{Numerical Stability.}
When computing $\mathbf{A}_m^{-1} \mathbf{b}_m$, we use the Sherman-Morrison formula for incremental updates rather than explicit matrix inversion, ensuring $O(d^2)$ complexity and numerical stability.

\subsubsection{Experimental Validation}

Crucially, we employ Dynamic Pareto Filtering in all BanditGPT experiments. For every incoming prompt $x$, the router constructs a local Pareto frontier based on contextual quality predictions $\hat{r}(x)$, restricting the search space to non-dominated models before the LinUCB selection policy is applied.

This mechanism is particularly critical in our 9-model experiments (Section~\ref{sec:results}), where it reduces the effective action space from 9 to $\approx$3 models per query, enabling efficient exploration without sacrificing optimality.

